\section{Counting Sundays --- Problem 19}
\label{prob:19}
\noindent\textbf{Source:} \url{https://github.com/KO0I/projectEuler/blob/reorganize/solutions-by-language/solutions/c/counting_sundays_19.c}

\noindent\textbf{Assessment:} Calendar simulation implemented.

\subsection*{Mathematical goal}
Count months whose first day is Sunday between January 1901 and December 2000, inclusive.

\subsection*{How the C code works}
The global state tracks year, month \verb|j|, weekday \verb|i|, day of month and \verb|score|. It starts on Tuesday, January 1, 1901. \verb|nextDay| cycles through weekdays, and \verb|nextMonth| advances to the first of the next month and increments \verb|score| when that day is Sunday. The main loop handles February and months with 30 or 31 days.

\subsection*{Correctness and limitations}
The test \verb|year%4==0| omits the full Gregorian century rule, but covers leap years in the specified 1901--2000 interval, including 2000. Formal function arguments are unused, and the stopping boundary deserves tests. The intended answer is 171.


\subsection*{Correction to historical C --- Problem 19}
{\color{red}\textbf{Correction:} The corrected calendar traversal uses the full Gregorian leap-year rule and counts Sundays at the start of each month only within the requested date interval.}
